% titel: Blatt 9, A4: Taylorentwicklung
% Teil II: der durchgerechnete Weg.

\subsection*{Gegeben und gesucht}
(a) Taylorreihe von $f(x)=\frac1{1+x^2}$ um $x_0=0$ auf $(-1,1)$; (b) daraus
arctan-Reihe durch gliedweise Integration; (c) $\arctan(0.5)$ per Python.

\subsection*{a) Taylorreihe von $\frac1{1+x^2}$}
Mit der geometrischen Reihe für $q=x^2$, $|x|<1$:
\begin{equation}
  f(x)=\frac1{1+x^2}=\sum_{k=0}^{\infty}(-1)^k x^{2k}
  =1-x^2+x^4-x^6+\cdots .
\end{equation}

\begin{ergebnis}
$\displaystyle f(x)=\sum_{k=0}^{\infty}(-1)^k x^{2k}$ für
$|x|<1$, mit $f^{(k)}(0)=0$ für ungerade $k$ und
$f^{(2k)}(0)=(-1)^k(2k)!$.
\end{ergebnis}

\subsection*{b) Taylorentwicklung von $\arctan$}
$\arctan$ ist Stammfunktion von $\frac1{1+t^2}$ mit $\arctan(0)=0$. Gliedweise
integrieren von $0$ bis $x$ (Satz 5.4.2):
\begin{align}
  \arctan(x)
  &= \int_0^x \sum_{k=0}^{\infty}(-1)^k t^{2k}\,\dd t
   = \sum_{k=0}^{\infty}(-1)^k \frac{x^{2k+1}}{2k+1} \\
  &= x-\frac{x^3}{3}+\frac{x^5}{5}-\frac{x^7}{7}+\cdots,\qquad |x|<1 .
\end{align}

\begin{ergebnis}
$\displaystyle\arctan(x)=\sum_{k=0}^{\infty}(-1)^k\frac{x^{2k+1}}{2k+1}$
für $|x|<1$.
\end{ergebnis}

\subsection*{c) $\arctan(0.5)$ mit Python}
\begin{verbatim}
import math
x = 0.5
N = 20
s = 0.0
for k in range(N):
    s += (-1)**k * x**(2*k+1) / (2*k+1)
print(s)                 # 0.463647609...
print(math.atan(0.5))    # Referenz
\end{verbatim}
Für $N=20$ erhält man $s\approx0.463647609$, in Übereinstimmung mit
$\mathrm{math.atan}(0.5)$. Da die Reihe alterniert, ist der Fehler durch das
erste weggelassene Glied $\tfrac{x^{2N+1}}{2N+1}$ schon sehr klein.

\begin{ergebnis}
$\arctan(0.5)\approx0.4636476$ (rad).
\end{ergebnis}

\begin{kontrolle}
\textbf{Referenz:} $\arctan(0.5)\approx0.4636$ rad $\approx26.6^\circ$ — passt,
da $\tan(30^\circ)=\frac1{\sqrt3}\approx0.577$. \textbf{Konvergenz:} $x=0.5<1$,
Reihe ist absolut konvergent, wenige Glieder genügen.
\end{kontrolle}

\begin{merksatz}
Potenzreihen aus der geometrischen Reihe über Substitution + gliedweise
Integration ableiten; die Integration ändert den Konvergenzradius nicht und
liefert die Stammfunktion als Reihe.
\end{merksatz}
